\chapter{Research and Existing Solutions}
\label{chap:research}
\section{Research Method / Approach}
We reviewed the approved proposal slides, panel action register, official mid-report template and rubric. We then read primary papers, original model cards and official product descriptions, and checked the saved scientific results and implementation. Product descriptions establish advertised features; they do not establish an independently audited evaluation protocol. The unsubmitted concept document is used as design reference rather than as an approved scope change.
\section{Existing Solutions / Products}
Wealthfront uses a questionnaire to recommend an allocation\cite{wealthfront_product}. PortfolioPilot offers portfolio analysis and personalized recommendations\cite{portfoliopilot_product}. Danelfin exposes technical, fundamental and sentiment subscores\cite{danelfin_product}. Magnifi offers conversational investment research\cite{magnifi_product}. Bloomberg and LSEG Workspace combine professional market data and news\cite{bloomberg_product,lseg_product}. StockTwits and TradingView provide social discussion or market-charting tools\cite{stocktwits_product,tradingview_product}. These products address different user tasks; Investify evaluates an inspectable historical research path rather than trying to replace a terminal or a licensed adviser.

Ruan and Jiang combine FinBERT-enhanced sentiment, XGBoost and SHAP in a published prediction pipeline\cite{ruan_jiang_2025}. This is close technical prior work, so combining these libraries is not itself our contribution. Li et al. identify training-window information leakage in LLM financial agents\cite{profit_mirage_2025}. Zhang et al.'s 2026 preprint changes one decision-time convention at a time and finds that inflation depends on the fault\cite{alpha_disappears_2026}; it is not peer reviewed. These studies motivate controlled timing comparisons and a declared null-result path.
\section{Comparison of Existing Solutions}
\begingroup\small
\begin{longtable}{p{.23\textwidth}p{.12\textwidth}p{.16\textwidth}p{.13\textwidth}p{.12\textwidth}p{.13\textwidth}}
\caption{Comparison on the five proposal criteria.}\label{tab:existing-solutions}\\\toprule Solution & Risk profile & Formal + informal sentiment & Explainable predictions & Leakage-safe evaluation & Verified natural language\\\midrule
Wealthfront & Risk quiz & Not specified & Allocation explanation & Not assessed & Not specified\\
PortfolioPilot & Risk preferences & Not specified as paired features & Portfolio recommendations & Not assessed & Not specified\\
Danelfin & Low-risk score; personal profile not established & Sentiment subscore; register separation unclear & Feature subscores & Not assessed & Not specified\\
Magnifi & Personal help & Not specified & Research chat & Not assessed & Not specified\\
Bloomberg / LSEG & Risk tools & News; other content varies & Analytics and research & Product-dependent; not assessed & Record-level guard not established\\
StockTwits / TradingView & Not established here & Social discussion / charts and news & Signals and charts & User-script dependent & Not established here\\
Published FinBERT + XGBoost + SHAP & Not central & News; no paired social experiment identified & SHAP attributions & Study-specific timing & No independent claim guard identified\\
Manual workflow & User judgement & User-selected sources & User reasoning & User-dependent & No automatic checker\\
Investify Mid-1 & Complete score + gate & News; social evaluated separately & All-feature SHAP context & Tested controls; Gate A partial & Template grammar + checked fallback\\
Investify planned & Same profile design & Integrated divergence ablation & SHAP + LLM prose & Broader timing/parity validation & Constrained claims, Mid-2\\
\bottomrule\end{longtable}
\endgroup
``Not specified'' means the inspected public source did not establish that criterion, not that the product lacks it. The manual row describes the workflow targeted by this project, not a measured user-study baseline.
\section{Limitations and Gaps}
The relevant gap is the combination of time-aware evaluation and a checkable explanation. Showing an indicator, a proprietary score or fluent prose does not by itself show that the inputs were available at decision time or that each explanation claim matches the computation. Conversely, checking a claim does not establish a useful predictor. The current project limitations are consolidated in Section~\ref{sec:limitations}.
\section{Proposed Solution}
The learner receives the model's probabilities, full evidence and a separate risk-policy outcome. A deterministic verifier checks known claims and retains the reason for rejection before emitting a checked replacement. The planned LLM translates the same record into accessible prose. Its design and evaluation are in Section~\ref{sec:llm-design}.
\subsection{Contribution and Evaluation}
\begin{enumerate}
\item \textbf{C1: Leakage-audited point-in-time harness with a permanent canary.} Implemented timing, scaling and label-end checks; evaluated by unit tests and fixed leakage arms. The centred Bollinger fault fires. Broader historical availability and the full proposal Gate A remain partial.
\item \textbf{C2: Deterministic claim verifier with retained rejection and checked fallback.} Implemented for the template grammar; evaluated using independent prose/JSON corruptions, omission tests and the live recovery path. Free-prose extraction is a Mid-2 deliverable.
\item \textbf{C3: Cross-register divergence features with a pre-registered ablation.} Planned. Evaluate a news-only model against news/social and divergence variants using shared chronological folds and settings, with the protocol fixed before running the comparison.
\end{enumerate}
As stated on proposal slide 8, no novelty is claimed in sentiment scoring, gradient boosting, SHAP-to-text or QLoRA, and no profit claim is made. The contributions concern controls and their evaluation, not the mere use of these components.
\section{Complex Computing Problem}
The requirement is to justify a stock-research outcome using observations that arrive on different clocks, preserve the information available at decision time, reconcile capacity with tolerance and keep every displayed claim consistent with the underlying computation. These requirements conflict when a fluent explanation or a personalized response would otherwise hide weak evidence.
\begin{longtable}{p{.26\textwidth}p{.63\textwidth}}
\caption{Seoul Accord characteristics applied to the project.}\label{tab:ccp}\\\toprule Characteristic & Project evidence and boundary\\\midrule
Conflicting requirements & Readability versus checkability; capacity versus tolerance; personalization versus privacy\\
No obvious solution & Timing and vendor revisions require an explicit information-availability model, not just a random train/test split\\
Depth of knowledge & Label overlap, calibration, attribution and claim verification require different mathematical models\\
Unfamiliar issues & Retroactive price adjustment, early closes and revised publication times create silent evaluation faults\\
Beyond routine practice & Team-built fault controls and record-based rejection tests complement libraries\\
Diverse stakeholders & Learner, operator and supervisor have different criteria, but the current operating roles are modest\\
Significant consequences & Unfounded certainty may mislead a non-expert interpreting a stock signal\\
Interdependence & Store, feature construction and evaluation are coupled; much of the pipeline is still sequential\\
Ill-defined requirements & Whether short-horizon predictive information exists must be discovered rather than assumed\\
\bottomrule\end{longtable}
This is a requirement-level argument adapted from concept Section 5. It does not treat every characteristic as equally strong.
\section{Technical Challenges}
The main challenges are timestamp precision, available-label boundaries, calibration collapse, social-domain transfer and consistent text/JSON claims. XGBoost\cite{xgboost_2016} and SHAP\cite{shap_2017,treeshap_paper} supply algorithms; the project must control how their inputs, outputs and claims are evaluated. Purged evaluation is motivated by overlapping financial labels\cite{deprado_2018}.
\section{Research Findings and Design Decisions}
The FNSPID precision audit led to a separate timestamped Alpaca/Benzinga experiment. Isotonic calibration collapsed the final classes; the subsequent bounded temperature fit produced near-uniform probabilities and motivated abstention. The full profile is used as a confidence gate rather than changing model direction. The social study motivates evaluating equity transfer before adding divergence. The current template verifier precedes any prompted LLM or QLoRA work\cite{qlora_2023,phi3_2024}.
\section*{Stock Market Background and Design Implications}
\addcontentsline{toc}{section}{Stock Market Background and Design Implications}
\label{sec:market-background}
\subsection*{Price Formation, Liquidity and Trading Sessions}
A share represents an ownership interest in a company. Secondary-market prices arise when buyers and sellers transact through trading venues. A bid is an offered buying price and an ask is an offered selling price; their difference is the spread. Market orders seek execution at available prices, while limit orders restrict the acceptable price. A daily close summarizes completed trading and is not a guarantee that a later order can be filled at that price\cite{sec_orders,sec_execution,sec_participants}. A model trained on daily bars cannot infer an intraday order book or assume unlimited liquidity.

Trading takes place on a calendar. US equity sessions contain holidays and early closes; five trading sessions can span more than five calendar days\cite{nyse_calendar}. An announcement after the close belongs to a different information set from one received before it. Our labels use observed price-session endpoints. The inherited fixed 16:00 New York news cutoff is recorded, including its early-close limitation.

\textbf{Design implication:} use an exchange calendar, preserve event and acquisition times and join completed news sessions backward to prices. Evaluate execution assumptions separately before any future trading claim. The present UI reports a historical date and horizon and has no order-execution path.
\subsection*{Price Drivers and the Short-Horizon Noise Problem}
Earnings and guidance can change expectations of future cash flows; valuation also depends on the return investors demand. Monetary policy and interest rates can affect discount rates, financing costs and alternatives to equities. Sector and index movements, institutional flows and liquidity can move a stock alongside its own news. These are reasons to avoid explaining a price change from a single headline. The project's eleven inputs contain price/volume summaries and headline sentiment; they are not a complete valuation or macroeconomic model.

Market efficiency treats prices as responding to available information\cite{fama_1970}. It does not prove that every market is always efficient or that no prediction can work. It does explain why widely observed short-horizon signals can have little remaining advantage once information, costs and timing are handled correctly. Positive historical performance requires a meaningful comparator and an out-of-sample protocol.

\textbf{Design implication:} compare with stratified-random and linear baselines, preserve weak findings and keep SHAP statements about model use separate from economic causation. A probability is a class estimate, not an expected return. The current no-signal result is a reason to abstain rather than to claim market insight.
\subsection*{The Five Technical Inputs}
Let $P_t$ be the completed closing price, $V_t$ traded share volume, and $\Delta P_t=P_t-P_{t-1}$. All windows below end at $t$; no centred window is allowed in the clean features. Indicator guides describe their usual interpretation\cite{fidelity_rsi,fidelity_macd,fidelity_obv}; the exact implemented formulas are in \nolinkurl{features/indicators.py}.
\begin{description}
\item[RSI(14)] With $g_t=\max(\Delta P_t,0)$ and $\ell_t=\max(-\Delta P_t,0)$, the implementation uses simple fourteen-session averages: $RS=\overline g_{14}/\overline\ell_{14}$ and $RSI=100-100/(1+RS)$. Completed all-gain/all-loss windows yield 100/0; a completed flat window yields 50. Warm-up remains missing.
\item[MACD histogram] $M_t=EMA_{12}(P)_t-EMA_{26}(P)_t$ and $H_t=M_t-EMA_9(M)_t$, where $EMA$ uses $\alpha=2/(n+1)$. This measures components of price momentum; its sign does not prescribe a trade.
\item[Distance from SMA200] $D_t=(P_t-SMA_{200,t})/SMA_{200,t}$. The long window means early rows do not have a complete value.
\item[Bollinger position] $B_t=(P_t-L_t)/(U_t-L_t)$, with $L_t=\overline P_{20}-2s_{20}$ and $U_t=\overline P_{20}+2s_{20}$. The sample standard deviation describes trailing variability. Values can lie outside $[0,1]$; a zero-width band is missing.
\item[OBV slope] $OBV_t=OBV_{t-1}+\operatorname{sign}(\Delta P_t)V_t$ and $O_t=(OBV_t-OBV_{t-14})/14$. This is a price-associated volume summary, not an observation of investors' intentions.
\end{description}
\textbf{Design implication:} expose the maintained definition and actual recorded value beside its SHAP effect. A value's sign and its attribution sign answer different questions. Direct perturbation checks verify that changing a future price cannot change an earlier clean feature.
\subsection*{The Six News Inputs, Absorption and Manipulation}
FinBERT assigns positive, neutral and negative headline probabilities\cite{finbert_paper,finbert_card}. For headline $i$, polarity is $q_i=p_i^{+}-p_i^{-}$. Let $n_t$ be the scored-headline count and $a_t$ the contributing session's average polarity. The implemented six news features are:
\begin{align*}
\text{mean\_polarity}_t &= a_t=\frac{1}{n_t}\sum_{i=1}^{n_t}q_i, & \text{mention\_volume}_t &= n_t,\\
\text{mom\_3d}_t &= \overline a_{3,t}-a_{t-3}, & \text{mom\_7d}_t &= \overline a_{7,t}-a_{t-7},\\
\text{volume\_z}_t &= \frac{n_t-\overline n_{20,t}}{s_{n,20,t}}, & \text{dispersion}_t &= s_{a,5,t}.
\end{align*}
Momentum allows shorter initial rolling windows but still needs the shifted prior value. Attention z-score needs at least five observations and a nonzero standard deviation; dispersion needs at least two. Dispersion is variation in session-average tone, not disagreement between individual headlines. Headline count is news attention, not traded share volume.

News can be anticipated, absorbed quickly or overwhelmed by other information. Tone can also reflect rumours, selective reporting or coordinated discussion. Neither a positive headline nor social agreement proves truth or price impact\cite{sec_ai_alert}. The paired social study therefore measures sentiment agreement on its frozen sample rather than investment performance. Cross-register divergence is planned and will be tested with shared evaluation settings.

\textbf{Design implication:} store the later creation/update timestamp for the new news snapshot, disclose coverage and missing features, and use a strict prior-session join. The future explainer must describe conflicting model contributions without inventing absent social values or claiming that a headline caused a return.
\subsection*{Volatility Regimes and Corporate Actions}
The distribution of returns and the usefulness of a signal may change between quiet and volatile periods. The same indicator value need not imply the same future outcome. Calibration fitted on one period may transfer poorly to a later refitted classifier, so mode and period identities remain visible.

Splits change share count and price units. NVIDIA's four-for-one 2021 split and ten-for-one 2024 split are documented by the issuer\cite{nvda_split_2021,nvda_split_2024}. The new historical prices use a current-vintage SIP split adjustment, including the vendor's action vintage, without cash-dividend adjustment. A discontinuity in raw prices must not be mistaken for an economic return; adjusted close and total return are different conventions.

\textbf{Design implication:} pin adjustment policy and data vintage, preserve issuer-action references and assess regimes and dividend treatment separately. Compare calibration on its declared fit/evaluation periods. The present model's uniform-confidence outcome supports no directional recommendation across either example profile.

\section*{Model Inventory and Adaptation}
\addcontentsline{toc}{section}{Model Inventory and Adaptation}
\label{sec:model-inventory}
\begingroup\small
\begin{longtable}{p{.23\textwidth}p{.29\textwidth}p{.37\textwidth}}
\caption{Checkpoints, tasks, choices, adaptation, licenses and status.}\label{tab:model-inventory}\\\toprule Model/checkpoint & Task, domain and inputs/outputs & Choice, adaptation, license and status\\\midrule
\nolinkurl{ProsusAI/finbert} & Financial headlines to positive/neutral/negative probabilities; polarity is positive minus negative & Chosen for financial text. Original weights, no project fine-tune. News inference completed. Current publisher weights card states CC BY-NC-SA 4.0; pinned-source license context retained.\\
\nolinkurl{ElKulako/cryptobert} & Crypto-social text to three-class sentiment; equity transfer evaluated & Original checkpoint, no project adaptation. MIT model card. Executed on frozen 24-post MSFT held-out set; macro-F1 0.2564.\\
\nolinkurl{cardiffnlp/twitter-roberta-base-sentiment-latest} & Twitter text to three-class sentiment & Alternative social domain. No project adaptation. CC BY 4.0 model card. Same 24 posts; macro-F1 0.6016.\\
XGBoost project model & Eleven price/news features to Buy/Hold/Sell scores and probabilities & Trained by the project; chronological validation. XGBoost software Apache 2.0. Retained frozen/refit model hashes in scientific receipts.\\
OpenAI GPT-6 assistant teacher & Saved record/context to controlled prose; six pilot examples & Teacher name comes from frozen plan. Exact serving revision unavailable. Proprietary service terms. JSON targets assembled deterministically; student not trained.\\
Microsoft Phi-3-mini & Planned local record/context-to-prose student & Original publisher family; exact project checkpoint not selected. MIT publisher model card. Prompted baseline Mid-2; QLoRA Iteration 4\cite{phi3_2024,qlora_2023}.\\
Mistral & Planned fallback local explainer & Exact fallback checkpoint unselected; Mistral-7B-Instruct-v0.3 is an Apache 2.0 publisher candidate. No executed fallback study.\\
\bottomrule\end{longtable}
\endgroup
The executed revisions are \nolinkurl{0bc2ac10dda6929bcb1f77ccc43c7ec5c424f031} (FinBERT), \nolinkurl{9e37c910fe87727cb842a9ac55c6388256fe0f15} (CryptoBERT), and \nolinkurl{3216a57f2a0d9c45a2e6c20157c20c49fb4bf9c7} (Cardiff). These are read from the research reference and paired study receipt. Model cards describe domain and licenses\cite{cryptobert_card,twitter_sentiment_card,phi3_card,mistral_card}; licenses are not inferred from the library used to load a model.

